\documentclass[10pt,a4paper]{amsart}
\usepackage[margin=19mm]{geometry}
\usepackage{amsmath,amssymb,mathtools}
\usepackage[scr=rsfs,cal=euler]{mathalfa}
\usepackage{xfrac}
\usepackage[unicode,hidelinks,bookmarks=false]{hyperref}
\newcommand{\ie}{i.e.\ }
\newcommand{\cf}{cf.\ }
\newcommand{\resp}{respectively\ }
\newcommand{\Subsection}{\subsection}
\providecommand{\nobreakdash}{\nobreak\hbox{-}}
\setlength{\emergencystretch}{2em}
\multlinegap0pt
\usepackage{mathtools}
\usepackage{amssymb}
\usepackage{bm}
\usepackage[shortlabels]{enumitem}
\usepackage{xcolor}
\usepackage{tikz}
\newtheorem{assumption}{Assumption}
\newtheorem{corollary}{Corollary}
\newtheorem{theorem}{Theorem}
\newcommand{\C}{\mathbb{C}}
\newcommand{\Z}{\mathbb{Z}}
\newcommand{\dist}{\operatorname{dist}}
\title{Spectral and Pseudospectral Approximation of Finite-Interaction-Range Operators in Doubling Metric Measure Spaces}
\author{Mattes Wittig}
\date{Source: arXiv:2608.03526v2 (CC BY 4.0)}
\begin{document}
\maketitle

\noindent\emph{Source and licence.} Spectral and Pseudospectral Approximation of Finite-Interaction-Range Operators in Doubling Metric Measure Spaces. Version arXiv:2608.03526v2 (CC BY 4.0), distributed by arXiv under the Creative Commons Attribution 4.0 International licence, http://creativecommons.org/licenses/by/4.0/, as recorded in the arXiv licence metadata for this exact version and shown on the abstract page https://arxiv.org/abs/2608.03526v2. Full licence notice: \texttt{licenses/arxiv-2608.03526-CC-BY-4.0.txt}.\par\smallskip

\noindent\textit{Editorial scope.} Counted unit: the article's theorem thm:pseudospectrum-nonnormal (source0.tex lines 1624--1645). The article's complete original proof of it is delivered with it (source0.tex lines 1647--1744, one \texttt{proof} environment of 98 lines). No mathematics is written, repaired or re-derived: the statement and the proof are the article's own lines, in the article's own notation, and no retained line is substituted or re-flowed.  Retained uncounted as the source-local context the proof needs: lines 567--573; lines 1536--1566.  Nothing was omitted from any retained span, so the package carries no omission record.  No retained line contains a citation, so the wrapper carries no bibliography and no bibliography entry.  No retained line contains a figure environment, an image-inclusion command or drawing code, so no image is delivered and the package declares no asset dependency.  Editorial formatting only, in addition: an AMS \texttt{amsart} preamble that restates the article's own declarations and macros used by the retained lines, namely the environments \texttt{assumption}, \texttt{corollary}, \texttt{theorem} on separate counters, together with the standard packages those lines need.  The retained environments print in this document as Assumption 1, Corollary 1, Theorem 1 in order of appearance, whereas the article prints the same items with the numbering of its own full document.  The complete native source of the article as distributed in the arXiv e-print for this exact version is bundled unchanged as \texttt{sources/206767/source0.tex}.
\begin{assumption}[Standing operator setting]\label{ass:operators}
From now on, let $\Gamma\subset X$ be
countable, uniformly discrete and relatively dense. Let
$H$ be a band operator on $\ell^2(\Gamma)$ with interaction range
$m>0$ and matrix bound $M>0$, in the sense introduced above. We denote by
$\gamma$ the corresponding interaction degree.
\end{assumption}

\begin{corollary}[Sampling the spectrum from local lower norms in the normal case]
\label{cor:normal-sampling-grid}
Let $\Gamma$ and $H$ satisfy Assumption \ref{ass:operators} and assume that $H$ is normal. 
For fixed $\delta>0$, choose $L$ so large that
\[
  \frac{2C_0}{L}<\delta.
\]
Set
\[
  h_L:=\frac{\sqrt2\,C_0}{L}
  \qquad\text{and}\qquad
  R>\|H\|+\frac{C_0}{L},
\]
and consider the finite grid
\[
  G_{L,R}
  :=\bigl\{a+ib\in\C:\ a,b\in h_L\Z,\ |a+ib|\le R\bigr\}.
\]
Define
\[
  \Lambda_L(\delta)
  :=\bigl\{\lambda\in G_{L,R}:\ \alpha_L(\lambda)<\delta\bigr\}.
\]
Then
\[
  d_H\bigl(\sigma(H),\Lambda_L(\delta)\bigr)<\delta.
\]
In particular, the spectrum of a normal finite--interaction--range operator can be
approximated arbitrarily well in Hausdorff distance by finitely many sample
points detected from local lower norms.
\end{corollary}

\begin{theorem}[Approximation of the pseudospectrum in the non-normal case]
\label{thm:pseudospectrum-nonnormal}
Let $H$ be as in Assumption \ref{ass:operators}, and fix $\varepsilon>0$ and $\delta>0$. For $L>L_0$, set
\[
h_L\coloneqq \frac{C_0}{L} \text{ and } R > \|H\|+\varepsilon+\frac{C_0}{L}.
\]
Let $G_{L,R}\subset\C$ be the grid with mesh size $h_L$ (as in Corollary \ref{cor:normal-sampling-grid}), and define
\[
\Lambda_\eta^{(h_L)}:=\{\mu\in G_{L,R}:\alpha_{L,comb}(\mu) < \eta\},
\qquad \eta > 0.
\]
Choose $L>L_0$ so large that
\[
h_L<\min \{ \frac{\delta}{4}, \frac{\varepsilon}{2} \}
\qquad\text{and}\qquad
d_H\bigl(\Lambda_{\varepsilon+2h_L}^{(h_L)},\Lambda_\varepsilon^{(h_L)}\bigr)<\frac{\delta}{2}.
\]
Then
\[
d_H\bigl(\sigma_\varepsilon(H),\Lambda_\varepsilon^{(h_L)}\bigr)<\delta.
\]
\end{theorem}

\begin{proof}
Since $\Lambda_\varepsilon^{(h_L)}\subset \sigma_\varepsilon(H)$, it suffices to show that
\[
\sup_{\lambda\in\sigma_\varepsilon(H)}
\dist\bigl(\lambda,\Lambda_\varepsilon^{(h_L)}\bigr)<\delta.
\]

The proof has two ingredients. First, one can indeed choose $L$ so large that
\[
d_H\bigl(\Lambda_{\varepsilon+2h_L}^{(h_L)},\Lambda_\varepsilon^{(h_L)}\bigr)<\frac{\delta}{2}.
\]
In the Globevnik case, the map
\[
\eta\longmapsto \sigma_\eta(H)
\]
is continuous with respect to the Hausdorff distance. So one can choose $L>L_0$ such that
\[
d_H\bigl(\sigma_{\varepsilon+2h_L}(H),\sigma_{\varepsilon-2h_L}(H)\bigr) < \frac{\delta}{4}.
\]
Let $\eta\in \Lambda_{\varepsilon+2h_L}^{(h_L)}$. Then $\eta\in \sigma_{\varepsilon+2h_L}(H)$. Hence we can choose
\[
\eta'\in \sigma_{\varepsilon-2h_L}(H)
\]
such that
\[
|\eta-\eta'|
\le d_H\bigl(\sigma_{\varepsilon+2h_L}(H),\sigma_{\varepsilon-2h_L}(H)\bigr) < \frac{\delta}{4}
\]
Since $G_L$ has mesh size $h_L$, there exists $\eta''\in G_L$ with
\[
|\eta'-\eta''|< h_L < \frac{\delta}{4}.
\]
Using the $1$-Lipschitz continuity of the combined lower norm, we obtain
\[
\nu_{comb}(H-\eta'')
\le \nu_{comb}(H-\eta')+|\eta'-\eta''|
< \varepsilon-2h_L+|\eta'-\eta''|
< \varepsilon-h_L.
\]
Hence
\[
\alpha_{L,comb}(\eta'')
\le \nu_{comb}(H-\eta'')+\frac{C_0}{L}
< \varepsilon-h_L+h_L
= \varepsilon,
\]
so that $\eta''\in \Lambda_\varepsilon^{(h_L)}$. Moreover,
\[
|\eta-\eta''|
\le |\eta-\eta'|+|\eta'-\eta''|
< \frac{\delta}{4}+h_L < \frac{\delta}{2}.
\]
\noindent
Thus we may choose $L$ so large that
\[
h_L<\frac{\delta}{4}
\qquad\text{and}\qquad
d_H\bigl(\Lambda_{\varepsilon+2h_L}^{(h_L)},\Lambda_\varepsilon^{(h_L)}\bigr)<\frac{\delta}{2}.
\]

Second, let $\lambda\in \sigma_\varepsilon(H)$. Since $G_L$ has mesh size $h_L$, there exists $\lambda'\in G_L$ such that
\[
|\lambda-\lambda'|<h_L.
\]
Therefore,
\[
\alpha_{L,comb}(\lambda')
\le \nu_{comb}(H-\lambda')+\frac{C_0}{L}
\le \nu_{comb}(H-\lambda)+|\lambda-\lambda'|+\frac{C_0}{L}
< \varepsilon+2h_L.
\]
Hence
\[
\lambda'\in \Lambda_{\varepsilon+2h_L}^{(h_L)}.
\]
It follows that
\[
\dist\bigl(\lambda,\Lambda_{\varepsilon+2h_L}^{(h_L)}\bigr)<h_L<\frac{\delta}{4}.
\]

Combining this with the choice of $L$, we obtain
\[
\dist\bigl(\lambda,\Lambda_\varepsilon^{(h_L)}\bigr)
\le
\dist\bigl(\lambda,\Lambda_{\varepsilon+2h_L}^{(h_L)}\bigr)
+
d_H\bigl(\Lambda_{\varepsilon+2h_L}^{(h_L)},\Lambda_\varepsilon^{(h_L)}\bigr)
<
\frac{\delta}{4}+\frac{\delta}{2}
<
\delta.
\]
Taking the supremum over all $\lambda\in \sigma_\varepsilon(H)$ yields
\[
d_H\bigl(\sigma_\varepsilon(H),\Lambda_\varepsilon^{(h_L)}\bigr)<\delta,
\]
as claimed.
\end{proof}

\end{document}
